\documentclass{article}
\usepackage{titlesec}
\usepackage{blindtext}
\usepackage{hyperref}
\usepackage[margin=0.7in]{geometry}
\usepackage[brazil]{babel}

\title{Projeto IA - AV1}
\date{}
\author{Paulo Henrico Rabelo - 2312652 \\ Levi Tabosa - 2224207}

\begin{document}
\maketitle
\tableofcontents

\section{Introdução}
Este relatório aborda técnicas fundamentais de Inteligência Artificial, com foco na implementação e análise de modelos de Regressão Linear. O trabalho foi dividido em dois estudos que exploram diferentes modelagens estatísticas, partindo de uma abordagem simples para uma análise de múltiplas variáveis.

O problema 1 consiste em analisar dados referentes aos pontos marcados por um \textit{quarterback} de futebol americano durante a temporada de 2004 da Liga Americana de Futebol. O tema é em torno da modelagem da relação matemática entre o número médio de jardas conquistadas por tentativa de passe (variável independente) e a pontuação final obtida pelo jogador (variável dependente).

O segundo problema foca no efeito da inspeção de raios X em circuitos integrados, baseado em dados de um artigo publicado na revista \textit{Electronic Packaging and Production} (2002). O objetivo é compreender como as doses de radiação absorvida (variável de resposta) são influenciadas conjuntamente pela corrente elétrica, medida em miliamperes, e pelo tempo de exposição, medido em minutos.

\subsection{O Objetivo das Análises de Regressão}
No aprendizado de máquina, modelos de regressão linear são fundamentais para calcular a relação entre variáveis. O intenção das análises é conseguir estimar a melhor reta ou hiperplano, assim minimizando o erro dessas previsões. A modelagem permite a compreensão da taxa de variação dos dados, como uma variável impacta a outra e fazer deduções preditivas para situações ainda não observados.

\subsection{Objetivo do Trabalho}
O objetivo do trabalho é desenvolver modelos de Regressão Linear Simples e Regressão Linear Múltipla. Seguindo as restrições da disciplina, a implementação foi feita em Python, utiliza a biblioteca \textit{Numpy}.

Visando melhorar o entendimento sobre estimação de coeficientes, e a capacidade de realizar estimativas pontuais e a conduzir avaliações de modelos utilizando métodos como o Coeficiente de Determinação ($R^{2}$ e $R^{2}$ ajustado), Erro Quadrático Médio (MSE) e Erro Absoluto Médio (MAE).

\end{document}
